\chapter{产品定位与白皮书边界}
\section{本章摘要}
本章说明 产品定位 在 BAI Work 技术体系中的位置。核心关键词包括 桌面智能体工作台、BAI/BAI Code 品牌统一、从副本独立后的边界治理、不杜撰未公开协议。本章内容以仓库当前实现为准, 以公开 BAI Code 安装资料为边界, 不把未公开的桌面本地 HTTP 协议写成既定官方契约。

\section{背景与目标}
本章聚焦产品定位。BAI Work 的设计起点不是重新发明一个孤立桌面客户端, 而是在 BAI Code 能力之上提供稳定的本地工作台。项目已经完成 BAI Work 独立品牌整理, 对外使用 BAI Work 名称, 对运行时使用 BAI Code 叙事, 对内部遗留命名采取兼容解释。这个选择的核心价值在于减少用户认知切换: 用户看到的是一个面向代码、调试、架构、文档、规划、评审、自动化和工具协同的桌面应用, 维护者看到的是一个可以逐步替换 runtime 后端而不破坏渲染层体验的壳层。

产品定位的目标可以概括为三点。第一, 保持桌面体验稳定, 让会话列表、消息流、工具结果、项目选择、模型设置和权限提示在不同运行时状态下有一致入口。第二, 尊重 BAI Code 公开文档的边界, 公开文档已经覆盖 baicode CLI、Python wheel 安装以及 OpenAI-compatible LLM 参数, 但没有公开桌面本地 HTTP service/session/event/permission/question 协议, 因此白皮书只能把当前实现写成 BAI Work 自有兼容桥和待确认协议缺口。第三, 将构建发布、凭据保护、EBAI 映射和 hooks 安全策略写成可执行工程约束, 便于后续团队按相同标准扩展。

历史兼容接口需要文档化治理。当前实现保留若干旧 schema 与 preload ABI, 用于稳定存量设置迁移。白皮书面向外部读者统一使用 BAI Work/BAI Code, 面向维护者时说明这些兼容接口的边界, 防止误删导致迁移失败。

EBAI 安装器会阻止 symlink 内容进入 rules skill, 复制目录时使用 verbatimSymlinks=false 并在发现 symlink 时抛错。这个细节很重要, 因为规则包来自外部仓库, 如果允许符号链接, 可能把用户不期待的本地路径打包到技能目录。

\section{仓库证据}
本白皮书基于当前仓库文件和公开 BAI Code 安装资料整理。关键证据包括: README.md 与 README.en.md 对产品定位, runtime 边界和官方契约缺口做了直接描述。 electron-builder.config.cjs 定义 extraResources, artifactName, mac/win/linux target 和 BAI\_WORK\_BUNDLE\_BAI\_CODE\_RUNTIME/BAI\_WORK\_BUNDLE\_BAI\_CODE\_OFFICIAL 的默认逻辑。 src/main/runtime/bai-work-adapter.ts 定义 DEFAULT\_BAI\_CODE\_COMMAND, DEFAULT\_BAI\_BASE\_URL, DEFAULT\_BAI\_MODEL, BAI-Code-Runtime, BAI-Code-Official, wheelhouse venv 准备流程和 /v1/* 兼容服务。 src/shared/app-settings-types.ts, src/shared/app-settings-provider.ts, src/shared/model-provider-presets.ts 描述 Provider, endpoint format, model profile, multimodal capability 和默认 BAI 预设。 src/main/services/ecc-mapping-service.ts 与对应测试描述 EBAI commands, agents, rules, hooks 的安装目标和启用约束。 resources/bai-code-runtime/README.md 记录 Intel runtime 使用 BAI Code 0.9.1 与 CPython 3.11.15 x86\_64。 resources/bai-code-official/README.md 记录 official wheels 来源, Python 3.10-3.13 要求, macOS arm64 和 Windows win\_amd64 wheelhouse 覆盖范围。 这些文件共同给出一个明确事实: 当前项目已经完成品牌切换、runtime 适配、wheelhouse 打包和 EBAI 映射的主要工程路径, 但仍保留若干历史 ABI 名称。文档在描述这些名称时使用“内部兼容名”或“设置迁移名”, 不把它们当作对外产品名称。

对外文档应避免两类错误。第一类是把未公开的 BAI Code 桌面 HTTP 协议当作既成事实, 这会误导二次开发者直接依赖不存在或未承诺的接口。第二类是把当前兼容桥误读成完整 BAI Code 服务实现, 这会掩盖 CLI bridge 在 approvals、questions、resume session、tool-level semantic events 上的能力差距。白皮书采用更保守的写法: 能从代码和资源中核对的写成当前实现, 只能由路线图推导的写成演进方向, 官方未公开的写成待确认。

运行时能力必须明确区分 available、unavailable 和 disabled。available 表示当前路径可用, unavailable 表示能力应该存在但受当前 runtime 或依赖限制不可用, disabled 表示产品在当前模式下不启用该能力。这个区分能帮助 UI 给出准确按钮状态和错误提示。

凭据来源有优先级。运行时会从 BAI credential settings、runtime override、selected provider profile、环境变量 BAI\_API\_KEY/BAI\_BASE\_URL/BAI\_MODEL 和默认值中解析有效配置。这个模型既允许设置页配置, 也允许调试时用环境变量覆盖, 但生产文档必须提醒不要把真实密钥写入脚本、截图或日志。

\section{关键设计}
在产品定位中, 设计约束来自产品目标、运行时现实和发布形态三方面。产品目标要求桌面壳对用户呈现完整工作台, 包括项目上下文、模型选择、长任务对话、侧边会话、文件引用和生成物管理。运行时现实要求兼容 baicode CLI 与 OpenAI-compatible endpoint, 同时不宣称拥有官方未公开的 session/event/permission/question 协议。发布形态要求 Intel、Apple Silicon、Windows 三类构建拥有不同的 runtime 资源策略, 但用户入口保持一致。

本章涉及的重点包括: 桌面智能体工作台、BAI/BAI Code 品牌统一、从副本独立后的边界治理、不杜撰未公开协议。这些重点之间存在明确依赖关系: 打包策略决定可用 runtime, runtime 能力决定本地 /v1/* 兼容层能否完整响应, 兼容层能力决定渲染层可以展示哪些状态和按钮, 安全策略决定哪些命令、hooks、文件访问和凭据路径必须默认收敛。BAI Work 的当前实现倾向于先稳定桌面壳, 再等待 BAI Code 官方契约成熟后替换兼容桥内部实现。

compatibility bridge 不是协议承诺, 而是过渡层。它在 /health 返回服务状态, 在 /v1/runtime/info 返回 provider、capabilities、sandboxMode、approvalPolicy 等信息, 在 /v1/threads 管理内存中的 thread 和 turn, 并通过 text/event-stream 输出 item\_updated、tool\_call\_started、tool\_call\_finished、assistant\_text\_delta、turn\_completed 等桌面事件。

OpenAI chat bridge 由 BAI\_WORK\_USE\_OPENAI\_CHAT\_BRIDGE=1 控制, 使用 chat/completions, stream=true, temperature=0, max\_tokens=1024。它能验证 OpenAI-compatible endpoint 路径和模型调用, 但不提供 BAI Code CLI 的本地工程执行语义。

\section{目标用户与使用场景}
软件工程师、技术负责人、企业内部工具团队、打包发布维护者 是本章最直接的读者。对这些读者而言, BAI Work 不只是一个聊天窗口, 而是本地工程活动的控制台。软件工程师关注是否能在项目目录内启动任务、读取上下文、运行必要命令并整理结果。发布维护者关注 runtime 资源是否随包交付, Python 版本是否满足 wheelhouse 要求, 构建产物命名是否稳定。安全人员关注 API key 是否只作为凭据处理, hooks 是否默认禁用, workspace 作用域是否明确。

典型使用场景包括本地代码修改、仓库评审、技术方案草拟、文档生成、调试辅助、自动化任务和跨渠道任务触发。不同场景对权限的要求不同: 文档草拟可以只需要模型调用和文件读取, 代码修改需要 workspace-write, shell 命令需要更高访问级别, hooks 则必须以受信任 workspace 为前提。白皮书把这些场景拆开描述, 目的是让用户在启用能力前知道能力边界, 让维护者在实现新功能时知道默认安全姿态。

历史兼容接口需要文档化治理。当前实现保留若干旧 schema 与 preload ABI, 用于稳定存量设置迁移。白皮书面向外部读者统一使用 BAI Work/BAI Code, 面向维护者时说明这些兼容接口的边界, 防止误删导致迁移失败。

EBAI 安装器会阻止 symlink 内容进入 rules skill, 复制目录时使用 verbatimSymlinks=false 并在发现 symlink 时抛错。这个细节很重要, 因为规则包来自外部仓库, 如果允许符号链接, 可能把用户不期待的本地路径打包到技能目录。

\section{实现说明}
当前实现中, 产品定位不是孤立模块, 而是跨主进程、渲染进程、共享设置和资源目录共同组成。主进程负责 runtime 探测、compatibility server、IPC handler、发布资源定位和本地文件写入。渲染进程负责工作台 UI、设置页、slash command 菜单、插件技能市场和用户交互。共享层负责类型、schema、Provider profile、endpoint format、默认值和路径规范。资源目录提供 bundled runtime、official wheelhouse、skills 和品牌资产。

这种分层有一个实用优点: 当 BAI Code 官方服务契约尚未公开时, 渲染层不必频繁改动。它仍然调用稳定的本地 /v1/* 桌面边界, 主进程可以在内部从 CLI bridge、OpenAI chat bridge 或将来的官方 service bridge 之间切换。这个边界也有代价: 兼容层必须诚实暴露 unavailable、disabled 或 unsupported 状态, 避免 UI 把尚未实现的 approvals、questions、resume-session 能力伪装成可用。

运行时能力必须明确区分 available、unavailable 和 disabled。available 表示当前路径可用, unavailable 表示能力应该存在但受当前 runtime 或依赖限制不可用, disabled 表示产品在当前模式下不启用该能力。这个区分能帮助 UI 给出准确按钮状态和错误提示。

凭据来源有优先级。运行时会从 BAI credential settings、runtime override、selected provider profile、环境变量 BAI\_API\_KEY/BAI\_BASE\_URL/BAI\_MODEL 和默认值中解析有效配置。这个模型既允许设置页配置, 也允许调试时用环境变量覆盖, 但生产文档必须提醒不要把真实密钥写入脚本、截图或日志。

\section{安全与合规考量}
安全策略采用最小惊讶原则。API key 不写入文档示例, 不写入日志, 不提交仓库, 示例统一使用 BAI\_API\_KEY=your-key。本地服务默认绑定 127.0.0.1, 避免暴露到局域网。运行时信息只报告 apiKeyConfigured 这样的布尔值, 不回显密钥。CLI 输出清理阶段会对当前配置中的 API key 做替换, 以降低错误展示凭据的风险。

EBAI hooks 是最高风险能力之一, 因为 hooks 本质上会执行本地 shell 命令。BAI Work 的安装器默认只在用户目录生成 disabled manifest, 并写入提示文档。启用 hooks 时必须提供 workspace, 结果写入该 workspace 下的 .bai/hooks.toml, 且 enabled = true 只存在于该受信任 workspace。这个策略不能消除 hooks 自身风险, 但能降低“安装即全局执行”的误用概率。

compatibility bridge 不是协议承诺, 而是过渡层。它在 /health 返回服务状态, 在 /v1/runtime/info 返回 provider、capabilities、sandboxMode、approvalPolicy 等信息, 在 /v1/threads 管理内存中的 thread 和 turn, 并通过 text/event-stream 输出 item\_updated、tool\_call\_started、tool\_call\_finished、assistant\_text\_delta、turn\_completed 等桌面事件。

OpenAI chat bridge 由 BAI\_WORK\_USE\_OPENAI\_CHAT\_BRIDGE=1 控制, 使用 chat/completions, stream=true, temperature=0, max\_tokens=1024。它能验证 OpenAI-compatible endpoint 路径和模型调用, 但不提供 BAI Code CLI 的本地工程执行语义。

\section{运维与发布影响}
产品定位会直接影响构建和现场部署。Mac Intel x64 可以依赖随包交付的 BAI-Code-Runtime, 该资源内含 baicode 入口和相对路径 Python/site-packages。Mac Apple Silicon arm64 与 Windows x64 走 official wheelhouse 路线, 目标机如果没有 Python 3.10-3.13, 应用无法创建本地 venv, 需要用户安装 Python 或在设置中指定已有 baicode 可执行文件。

发布侧需要关注 electron-builder 的条件打包逻辑。BAI\_WORK\_BUNDLE\_BAI\_CODE\_RUNTIME 默认只在 darwin x64 打开, BAI\_WORK\_BUNDLE\_BAI\_CODE\_OFFICIAL 默认在 darwin arm64 和 win32 x64 打开。artifactName 采用 BAI-Work-{version}-{os}-{arch}.{ext}, update channel 只接受 stable 或 frontier。白皮书不把这些变量写成用户必须配置的运行条件, 而是写成发布工程可调整的构建参数。

历史兼容接口需要文档化治理。当前实现保留若干旧 schema 与 preload ABI, 用于稳定存量设置迁移。白皮书面向外部读者统一使用 BAI Work/BAI Code, 面向维护者时说明这些兼容接口的边界, 防止误删导致迁移失败。

EBAI 安装器会阻止 symlink 内容进入 rules skill, 复制目录时使用 verbatimSymlinks=false 并在发现 symlink 时抛错。这个细节很重要, 因为规则包来自外部仓库, 如果允许符号链接, 可能把用户不期待的本地路径打包到技能目录。

\section{边界、限制与后续演进}
当前最大限制仍是官方桌面服务契约未公开。BAI Work 可以调用 baicode CLI, 可以用 OpenAI-compatible chat completions 做有限 bridge, 可以为渲染层模拟 thread/turn/event 结构, 但不能声称已经获得官方 session/event/permission/question 协议。这个限制影响恢复会话、细粒度权限审批、问题回传、真实工具事件映射、长期存储和多客户端同步。

后续演进应遵循两条线。第一条是产品线: 统一品牌语言、简化设置项、提高错误可读性、完善缺少 Python 或缺少 API key 时的引导。第二条是工程线: 在官方契约发布后替换 compatibility bridge 内部实现, 保留 renderer 边界, 增加协议测试, 增加跨平台安装 smoke test, 补齐 hooks 审计和 Provider 能力探测。

运行时能力必须明确区分 available、unavailable 和 disabled。available 表示当前路径可用, unavailable 表示能力应该存在但受当前 runtime 或依赖限制不可用, disabled 表示产品在当前模式下不启用该能力。这个区分能帮助 UI 给出准确按钮状态和错误提示。

凭据来源有优先级。运行时会从 BAI credential settings、runtime override、selected provider profile、环境变量 BAI\_API\_KEY/BAI\_BASE\_URL/BAI\_MODEL 和默认值中解析有效配置。这个模型既允许设置页配置, 也允许调试时用环境变量覆盖, 但生产文档必须提醒不要把真实密钥写入脚本、截图或日志。

\section{工程判定清单}
\begin{itemize}
  \item 对外文案必须使用 BAI Work 和 BAI Code, 历史兼容名称只作为兼容名解释。
  \item 所有示例凭据必须使用占位符, 例如 BAI\_API\_KEY=your-key。
  \item 涉及 session、event、permission、question 的官方协议时必须写明“待官方公开契约确认”。
  \item 涉及 hooks 时必须写明默认 disabled, 启用必须 trusted workspace scoped。
  \item 涉及 Apple Silicon 或 Windows runtime 时必须写明 Python 3.10-3.13 要求。
\end{itemize}

\section{专项展开: 品牌治理}
用户不会关心内部迁移过程, 只会关心桌面应用是否稳定、能力边界是否清楚、失败提示是否能指导下一步。BAI Work 的产品定位应围绕“本地工程工作台”展开, 以项目为组织单位, 以 BAI Code 为执行核心, 以权限和配置为治理边界。历史来源需要写入白皮书, 但不能成为产品叙事中心。正确做法是说明它完成 BAI Work 独立品牌整理, 已切换对外名称、包名、图标、发布产物和 README 叙事, 同时仍保留若干内部兼容名。这样既不隐瞒迁移现实, 也不让用户误以为当前产品仍依附旧品牌。

从维护视角看, 产品定位需要持续回答四个问题。第一, 用户看到的能力是否已经由当前代码或资源支撑。第二, 失败时是否给出可执行的恢复路径。第三, 默认配置是否保护本地文件、凭据和工作区。第四, 新增能力是否会扩大打包、启动、权限或协议依赖。BAI Work 当前处在品牌切换和 runtime 适配并行推进阶段, 因此文档必须把“当前实现”“兼容层”“待确认官方契约”和“路线图”分开写。只有这样, 后续维护者才能在不破坏用户体验的前提下替换内部实现。

对用户而言, 产品定位最重要的是可预测性。应用启动后应该能说明 runtime 是否可用, baicode 是否可探测, API key 是否已配置, Provider base URL 是否明确, workspace 是否已选择, 以及当前权限是否允许文件修改和命令执行。对发布团队而言, 可预测性意味着构建参数、资源目录、产物名称、签名状态和 update channel 都可复现。对安全团队而言, 可预测性意味着 hooks 不会在用户不知情的情况下执行, 凭据不会出现在日志或文档中, 本地服务不会监听外部网卡。

在实现层面, 产品定位还应避免把历史兼容当作永久架构。兼容名可以保留, 因为它们保护旧设置和旧 ABI, 但新增文档、错误提示和用户界面应尽量使用 BAI Work/BAI Code 术语。compatibility bridge 可以保留, 因为它保护 renderer 稳定边界, 但它内部的模拟 thread store、CLI 输出整理和 unavailable route 必须随着官方契约成熟而替换。EBAI 映射可以保留, 因为它帮助用户迁移已有规则资产, 但 hooks 必须继续采用默认禁用和受信任 workspace 作用域。

验收 产品定位 的方式也应具体。文档验收看是否列明路径、默认值、限制和恢复步骤。代码验收看是否有单元测试覆盖成功路径和失败路径。发布验收看目标平台是否包含预期资源, Intel 包是否包含 BAI-Code-Runtime, Apple Silicon 和 Windows 包是否包含 BAI-Code-Official, wheelhouse 是否覆盖 cp310 到 cp313。安全验收看是否没有真实密钥, 是否没有把 hooks 写成默认启用, 是否没有把未公开协议写成官方承诺。

这些要求共同形成 BAI Work 的文档纪律: 宁可把不确定性写清楚, 也不为了完整感补齐不存在的协议细节。宁可把跨平台差异写复杂一点, 也不让 Apple Silicon 或 Windows 用户误以为无需 Python。宁可把 hooks 风险重复说明, 也不让用户把外部命令包当作普通提示词资产。宁可承认内部命名仍在迁移, 也不为了品牌统一在代码层做高风险重命名。

\section{维护者补充说明}

维护者在处理产品定位相关变更时, 应先确认变更是否直接服务用户请求或当前路线图。BAI Work 已经有较多跨层能力, 包括桌面 UI、运行时桥接、Provider 设置、技能市场、EBAI 安装、发布脚本和多平台资源。任何看似简单的改动, 都可能影响打包产物、首启流程、设置迁移或安全默认值。比如修改默认模型会影响 Provider preset、runtime config、composer 默认值和文档示例。修改数据目录会影响 venv 缓存、工作区默认路径和用户已有状态。修改 hooks 事件映射会影响外部命令执行时机。

因此, 白皮书把产品定位写成可维护对象, 而不是一次性说明文字。每个章节都应帮助读者回答“当前版本如何工作”“为什么这样设计”“哪里不能承诺”“如何验证”“下一步如何演进”。这种写法会比营销式介绍更长, 但能减少后续误解。特别是在 BAI Code 官方桌面 service 契约尚未公开的阶段, 文档越需要克制: 可以写 BAI Work 当前有 compatibility bridge, 可以写 renderer 依赖 /v1/* 桌面边界, 可以写某些 route 返回 unsupported 或 unavailable, 但不能写官方已经承诺这些 HTTP path。

在团队协作中, 产品定位还需要形成检查清单。提交前检查是否修改了写入范围之外的文件, 是否引入真实凭据, 是否把示例 key 写成真实格式, 是否使用了不适合 XeLaTeX 的字符, 是否让 README 的编译命令与实际目录不一致, 是否让 references.bib 引用不存在来源。发布前检查包内是否存在预期资源, Intel 包是否有 BAI-Code-Runtime, arm64 和 Windows 包是否有 BAI-Code-Official, Windows 安装器是否保留用户 app data, macOS 包是否使用正确图标和 entitlements。

从用户教育角度看, 产品定位应使用直接语言。用户不需要知道每一个内部函数名, 但需要知道配置 API key 后会发生什么, 缺少 Python 时该做什么, 开启 hooks 前要信任什么, 自定义 Provider 需要满足什么接口形态, 为什么某些权限或问题功能当前不可用。白皮书把这些解释写入正文, 是为了让产品、开发、测试、安全和发布团队在同一份材料上形成共识。

\section{验收与交付说明}
产品定位的交付验收应同时覆盖文档、代码、打包和安全四个层面。文档层面要求章节描述与仓库事实一致, 不出现聊天痕迹, 不出现真实凭据, 不把待确认协议写成官方能力。代码层面要求相关实现能在仓库中找到对应文件, 例如运行时适配应能追溯到主进程 adapter, Provider 能力应能追溯到共享设置和 preset, EBAI 行为应能追溯到安装服务和测试。打包层面要求目标平台资源与构建配置一致, 不把 Intel runtime 和 official wheelhouse 的适用平台混写。安全层面要求 hooks、shell、workspace、API key 和本地端口都有明确默认值和风险说明。

交付给新开发者时, 产品定位章节应帮助其快速判断自己要改哪里。如果要调整用户界面, 先看渲染层和共享文案。如果要调整本地 runtime, 先看主进程 adapter 和运行时设置。如果要调整模型列表, 先看 Provider preset 和能力 profile。如果要调整 EBAI 安装, 先看服务实现和测试。如果要调整发布产物, 先看 electron-builder 配置和 release 脚本。这样的路径说明能减少盲目搜索, 也能减少跨目录误改。

交付给产品和安全评审时, 产品定位章节应说明取舍原因。BAI Work 选择本地 compatibility bridge, 是为了在官方桌面服务契约尚未公开时保持桌面体验连续。BAI Work 选择 wheelhouse 加目标机 Python, 是为了在 Apple Silicon 和 Windows 上复用官方 BAI Code wheel 安装模型。BAI Work 选择 hooks 默认禁用, 是因为外部 hooks 会执行本地命令。BAI Work 保留部分历史内部命名, 是因为设置迁移和 ABI 稳定比一次性重命名更重要。

交付后的维护还需要周期性复查。每当 BAI Code 官方文档更新、wheel 版本变化、Provider 模型列表变化、Electron 或 electron-builder 升级、Windows 安装器策略变化、macOS 签名要求变化、EBAI 上游结构变化时, 白皮书都应同步更新。否则文档会从工程证据变成历史快照。维护者应优先更新事实性段落、命令、路径、默认值和已知限制, 再考虑扩展路线图。

\section{术语一致性说明}
产品定位章节使用的术语应与整份白皮书保持一致。对外统一写 BAI Work、BAI Code、BAI provider、BAI-Code-Runtime、BAI-Code-Official。涉及历史实现时, 可以写内部兼容名、设置迁移名或旧 ABI 名, 但必须说明这些名称不代表当前对外品牌。涉及官方能力时, 只写公开文档已经覆盖的 baicode CLI、Python wheel 安装和 OpenAI-compatible LLM 参数。涉及桌面本地 service、session、event、permission、question 时, 必须写成 BAI Work 当前兼容桥或待官方公开契约确认。

这种术语纪律看似细碎, 实际上直接影响项目可信度。用户会根据文档判断应用是否成熟, 集成方会根据文档判断接口是否稳定, 安全评审会根据文档判断风险是否被承认。如果同一能力在不同章节被写成不同名称, 或把过渡实现写成官方契约, 后续沟通成本会明显升高。因此每次更新产品定位相关内容时, 都应同时检查标题、表述、路径、默认值和限制说明是否仍然一致。

\section{更新原则}
白皮书后续更新应坚持证据优先。新增事实必须能追溯到代码、配置、资源目录、测试或公开文档。无法验证的判断只能写成假设、风险或路线图, 不能写成当前能力。这个原则适用于所有章节, 也适用于 README 中的编译命令和统计口径。

\section{小结}
从产品定位视角看, BAI Work 的短期任务是把已经完成的跨平台构建和 BAI Code 接入能力稳定交付给用户, 同时把不确定协议边界暴露给维护者。长期任务是在 BAI Code 官方桌面服务契约明确后, 用正式 service/session/event/permission/question 实现替换当前 bridge 内部逻辑, 但保留桌面壳、设置模型和用户工作流的连续性。
